\chapter{Divorce and Break-Up}

\section*{Definition and Overview}
Divorce and relationship break-ups are major life events that involve significant emotional, practical, and sometimes financial upheaval. The end of a relationship can trigger grief, stress, anxiety, and depression, and it affects every aspect of life, including children, housing, finances, and social networks. Most people go through a difficult period of adjustment and eventually recover, but the process is easier with support, self-care, and practical planning. This chapter explains the common experiences after a break-up, when to seek help, and how to support children and yourself.

\section*{Epidemiology}
Divorce is common in Australia, with around a third of marriages ending in divorce, and many more de facto relationships ending. Separation and divorce are major risk factors for psychological distress, depression, and anxiety, and are associated with increased use of health services. Children of separating parents experience significant adjustment, although most adapt well with supportive parenting. Around a quarter of Australian children experience their parents' separation, making this an important issue for family health.

\section*{Aetiology and Pathophysiology}
\begin{itemize}
    \item \textbf{Grief:} the loss of the relationship triggers grief, even when the break-up was wanted
    \item \textbf{Stress:} separation activates the body's stress response, affecting sleep, appetite, and health
    \item \textbf{Loss of identity:} roles, routines, and social networks change
    \item \textbf{Practical pressures:} finances, housing, and parenting arrangements add stress
    \item \textbf{Impact on children:} children respond to conflict, change, and the wellbeing of their parents
    \item \textbf{Recovery:} adjustment occurs over months to years, influenced by support and coping
\end{itemize}

\section*{Clinical Features}
\begin{itemize}
    \item Sadness, grief, anger, and feelings of failure or guilt
    \item Anxiety about the future, finances, and children
    \item Sleep disturbance, poor appetite, and fatigue
    \item Difficulty concentrating and reduced work performance
    \item Withdrawal from friends and activities
    \item Depression and anxiety that can persist beyond the expected adjustment period
    \item Physical symptoms such as headaches and stomach upsets
\end{itemize}

\section*{Differential Diagnosis}
\begin{itemize}
    \item Normal grief and adjustment versus depression requiring treatment
    \item Anxiety disorders triggered by the separation
    \item Alcohol or substance use as a coping strategy
    \item Difficulty with children's adjustment, which is distinct from the parent's distress
\end{itemize}

\section*{When to Seek Medical Care}
Seek help if low mood, anxiety, or distress persists, affects daily function, or leads to harmful coping such as heavy drinking. Parents should seek support for children who show persistent behavioural or emotional changes. A GP can provide support, counselling referrals, and treatment.

\textbf{Call 000 (or local emergency services) for:}
\begin{itemize}
    \item Thoughts of self-harm or suicide
    \item Concerns for the safety of a child or a family member
    \item Severe distress, agitation, or inability to cope with danger
    \item Domestic violence or immediate risk of harm
\end{itemize}
\textbf{Support lines:} Lifeline (13 11 14), 1800RESPECT (1800 737 732), and MensLine Australia (1300 789 978).

\section*{Diagnosis and Investigations}
\begin{itemize}
    \item \textbf{Psychological assessment:} mood, anxiety, coping, and functioning
    \item \textbf{Risk assessment:} suicide risk, substance use, and family violence
    \item \textbf{Practical assessment:} financial, housing, and parenting pressures
    \item \textbf{Children's wellbeing:} review of children's adjustment and needs
    \item \textbf{Referral:} to counselling, financial counselling, or legal services as needed
\end{itemize}

\section*{Management}
\begin{itemize}
    \item \textbf{Psychological support:} counselling, including grief and separation counselling
    \item \textbf{Peer support:} friends, family, and support groups
    \item \textbf{Self-care:} sleep, nutrition, exercise, and time for recovery
    \item \textbf{Practical planning:} financial counselling, legal advice, and housing support
    \item \textbf{Co-parenting support:} mediation and parenting plans that put children first
    \item \textbf{Support for children:} age-appropriate communication and professional help when needed
    \item \textbf{Treatment:} therapy and, where appropriate, medication for depression or anxiety
\end{itemize}

\section*{Complications}
\begin{itemize}
    \item Depression and anxiety disorders
    \item Alcohol and drug misuse
    \item Financial hardship and housing instability
    \item Conflict with an ex-partner, including family violence
    \item Adjustment difficulties in children
    \item Social isolation and loss of support networks
\end{itemize}

\section*{Prognosis and Outcomes}
Most people adjust to separation and divorce over one to two years, and many report personal growth and improved wellbeing afterwards. Support, healthy coping, and cooperative co-parenting are associated with better outcomes for both adults and children. Depression and anxiety are treatable, and most people who seek help recover fully. With time and support, most families adapt to their new circumstances.

\section*{Prevention and Self-Care}
\begin{itemize}
    \item Allow yourself to grieve and process the loss
    \item Keep routines of sleep, meals, and activity
    \item Stay connected with supportive people
    \item Limit alcohol, which worsens distress
    \item Seek counselling early rather than coping alone
    \item Talk to children openly and reassure them about their safety
    \item Use mediation and legal services to manage practical matters calmly
    \item Look after your physical health during this stressful time
\end{itemize}

\section*{Sources and Further Reading}
The following reputable sources provide further information for healthcare students and patients seeking reliable, current advice about divorce and break-up.

\begin{itemize}
    \item Relationships Australia. \textit{Support for separation and divorce.} https://www.relationships.org.au/
    \item Beyond Blue. \textit{Coping with separation and relationship breakdown.} https://www.beyondblue.org.au/
    \item Lifeline Australia. \textit{Support and crisis services.} https://www.lifeline.org.au/
    \item Raising Children Network. \textit{Separation, divorce, and children.} https://raisingchildren.net.au/
    \item Federal Circuit and Family Court of Australia. \textit{Information about family law and separation.} https://www.fcfcoa.gov.au/
    \item NHS. \textit{Managing stress and mental health after separation.} https://www.nhs.uk/mental-health/
\end{itemize}
